% 03_nlft_geometry.tex
\section{Nonlinear Fourier geometry of phase finding}
\label{sec:nlft}

The equivalence of generalized QSP with the $\SU(2)$ nonlinear
Fourier transform was established
in~\cite{AlexisMnatsakanyanThiele2024,AlexisBeckerEtAl2026,Laneve2025},
and the layer-stripping inverse was made
algorithmic in~\cite{NiSarkarYingLin2025}.  What the stereographic
chart adds is a geometric reading: the transfer matrix and the
signal operator $O_z$ are the same object on $\CP$, and the
backward direction of Theorem~\ref{thm:gen-qsp} is layer stripping
with the Schur parameters read as points of $\CP$.

\subsection{The transfer matrix is the signal operator}
\label{subsec:transfer}

The $\SU(2)$ NLFT of a sequence $(\gamma_0,\dots,\gamma_{d-1})$ is
the ordered product
\begin{equation}
	\widehat\gamma(z)=\prod_{k=0}^{d-1}
	\frac{1}{\sqrt{1+|\gamma_k|^2}}
	\begin{pmatrix} 1 & \gamma_k z^k\\ -\bar\gamma_k z^{-k} & 1
	\end{pmatrix},\qquad z\in\T,
	\label{eq:nlft}
\end{equation}
whose $z$-powers absorb into a diagonal signal so that the product
telescopes into a QSP sequence with $z$-independent
factors~\cite{Laneve2025}.  Each factor is the stereographic signal
operator up to a fixed rotation:
\begin{equation}
	\boxed{\ \begin{pmatrix} 1 & \gamma\\ -\bar\gamma & 1
	\end{pmatrix}=O_\gamma\,(i\PauliY),\qquad
	i\PauliY=\begin{pmatrix}0&1\\-1&0\end{pmatrix}\in\SU(2).\ }
	\label{eq:transfer-identity}
\end{equation}
Both sides have determinant $1+|\gamma|^2$, and both encode the
point $\gamma\in\CP$ in a matrix whose determinant is the squared
Fubini--Study norm of $(\gamma:1)$.  The normalization
$\sqrt{1+|\gamma|^2}$ common to the NLFT and to stereographic QSP
is this norm, which is why it appears in both.  Interleaved with
$\SU(2)$ controls, the fixed factor $i\PauliY$ absorbs into the
controls at each step, so the two product structures coincide.

\subsection{Layer stripping and the Schur parameter}
\label{subsec:stripping}

\begin{proposition}[Backward direction is layer stripping]
	\label{prop:stripping}
	The backward direction of Theorem~\ref{thm:gen-qsp} is the
	inverse-NLFT layer-stripping algorithm
	of~\cite{NiSarkarYingLin2025,AlexisBeckerEtAl2026} expressed in
	the stereographic coordinate.  The control $G_d$ extracted at
	the outer layer satisfies $(a_d,b_d)\propto(p_d,q_d)$, where
	$p_d,q_d$ are the leading coefficients of the reduced pair; the
	divisibility $1+z\bar z\mid P'$ that lowers the bidegree is the
	stereographic form of the NLFT condition that the residual datum
	be divisible by $z$; and after stripping,
	$|P|^2+|Q|^2=(1+|z|^2)^{d-1}$, which on $\T$ is the standard
	NLFT normalization up to the scale $2^{-(d-1)/2}$.
\end{proposition}

The Schur parameter at each layer is the homogeneous coordinate
$[p_d:q_d]\in\CP$.  Its $\SU(2)$ lift is the Hopf map
\begin{equation}
	(a_d,b_d)=\frac{(p_d,q_d)}{\sqrt{|p_d|^2+|q_d|^2}},
	\label{eq:hopf}
\end{equation}
so the inhomogeneous coordinate $\gamma_d=q_d/p_d$ is the classical
Schur parameter.  Working with the homogeneous pair keeps the case
$p_d=0$ regular: the point $[0:q_d]$ is an ordinary point of $\CP$,
whereas the classical NLFT excludes it by requiring $a^*(0)>0$.
The stereographic picture treats all Schur parameters uniformly as
points of $\CP$, and this is the only structural difference from
the standard layer-stripping recursion.  The proof of
Proposition~\ref{prop:stripping} is the backward direction of
Theorem~\ref{thm:gen-qsp}, read against the NLFT dictionary; we
give the correspondence in Appendix~\ref{app:thm-gen}.

\subsection{The Riccati flow}
\label{subsec:riccati}

\begin{remark}[QSP as a discrete Riccati flow]
	\label{rem:riccati}
	The M\"obius action of Sec.~\ref{subsec:encoding} makes each
	layer a fractional-linear step on $\CP$.  Tracking the encoded
	coordinate through the product of Theorem~\ref{thm:gen-qsp}
	gives the discrete Riccati iteration
	$\zeta_k=G_k\cdot O_z\cdot\zeta_{k-1}$, which for real signal
	and $\PauliZ$-phase controls reads
	$\zeta_k=e^{2i\phi_k}(x\zeta_{k-1}-1)/(\zeta_{k-1}+x)$.  Its
	continuum limit is the scalar Riccati equation
	$i\dot\zeta=\nu(1-\zeta^2)-2\phi(t)\zeta$ obtained by projecting
	the Zakharov--Shabat system
	$i\dot\psi=(\nu\PauliX+\phi(t)\PauliZ)\psi$ to
	$\zeta=\psi_2/\psi_1$, the Ablowitz--Kaup--Newell--Segur
	scattering problem~\cite{AblowitzKaupNewellSegur1974,FaddeevTakhtajan1987}
	with $\phi$ the potential and $\nu$ the spectral parameter.  The
	Schur parameters are the Verblunsky coefficients of the
	orthogonal polynomials on the unit circle associated with the
	iteration~\cite{Simon2005OPUC}, and the reflectionless case
	connects to the canonical systems of
	Ref.~\cite{BessonovLukicYuditskii2024}.  This correspondence is
	a \emph{description} of the QSP iteration, not an alternative
	algorithm: it situates phase finding within inverse scattering
	and explains why the layer-stripping and NLFT machinery apply,
	but the protocol itself remains the discrete
	iteration~\cite{TaoThieleTsai,NiSarkarYingLin2025}.
\end{remark}
